\documentclass[11pt]{article}
% Shared print layout. The web's provisional design and MathJax are independent.
\usepackage[utf8]{inputenc}
\usepackage[T1]{fontenc}
\usepackage[spanish,es-nodecimaldot]{babel}
\usepackage{lmodern}
\usepackage{amsmath,amssymb}
\usepackage{graphicx}
\usepackage{geometry}
\geometry{a4paper,left=20mm,right=20mm,top=25mm,bottom=22mm,headheight=16pt,headsep=8mm}
\usepackage{fancyhdr,lastpage}
\usepackage{enumitem}
\usepackage{xcolor}
\usepackage{pgfplots}
\pgfplotsset{compat=1.18}
\usepgfplotslibrary{fillbetween}
\usepackage[hidelinks,pdfencoding=auto]{hyperref}
\hypersetup{pdflang=es-ES,pdfauthor={},pdfsubject={Matemáticas Aplicadas a las Ciencias Sociales II}}
\pagestyle{fancy}
\fancyhf{}
\fancyhead[L]{\small Matemáticas Aplicadas a las Ciencias Sociales II}
\fancyhead[R]{\small 2.\textsuperscript{o} Bachillerato}
\fancyfoot[L]{\footnotesize Probabilidad II | Edición revisada 1.0}
\fancyfoot[R]{\footnotesize Página \thepage\ de \pageref{LastPage}}
\renewcommand{\headrulewidth}{0.4pt}
\setlength{\parindent}{0pt}
\setlength{\parskip}{0.5em}
\setlist{itemsep=0.3em,topsep=0.3em}
\newcommand{\apartado}[1]{\subsection*{#1}}
\pgfplotsset{
  revision/.style={
    width=0.94\linewidth,height=7.5cm,
    axis lines=middle,grid=major,grid style={gray!18},
    tick label style={font=\small},label style={font=\small},
    xlabel={$x$},ylabel={$y$},samples=160,
    legend style={at={(0.5,-0.16)},anchor=north,draw=none,font=\small},
    clip=true
  }
}

\hypersetup{pdftitle={Probabilidad II - Examen revisado}}
\begin{document}
\begin{minipage}{0.17\linewidth}\includegraphics[width=0.8\linewidth]{logo-al-andalus.png}\end{minipage}\hfill
\begin{minipage}{0.79\linewidth}{\Large\bfseries Probabilidad II}\\[0.4em]
Duración: \textbf{55 minutos}. Puntuación: \textbf{10 puntos}.\end{minipage}

Apellidos y nombre: \hrulefill

Fecha: \hrulefill

Responde razonadamente y numera tus páginas. Se permite calculadora. Cada ejercicio incluye 0,5 puntos de expresión de procesos y métodos, sin añadirlos a su máximo. Resuelve el ejercicio obligatorio (3 puntos) y solo dos de los tres opcionales (3,5 puntos cada uno).




\section*{Bloque obligatorio}
\subsection*{Ejercicio 1 \hfill 3 puntos} Sean dos sucesos \( A \) y \( B \) tales que \( P(A) = 0.4 \), \( P(B) = 0.5 \), y \( P(A \cap B^c) = 0.25 \).

\begin{enumerate}[label=(\alph*)]

\item (1 punto)  Calcule la probabilidad de que ocurra \( A \) y ocurra \( B \).



\item (1 punto)  Calcule la probabilidad de que no ocurra \( A \) pero sí ocurra \( B \), es decir, \( P(A^c \cap B) \).



\item (0,5 puntos)  Calcule la probabilidad de que ocurra \( A \) sabiendo que ha ocurrido \( B \), es decir, \( P(A \mid B) \).



\item (0,5 puntos)  \textquestiondown{}Son independientes \( A \) y \( B \)?



\end{enumerate}



\section*{Bloque opcional: elige dos de tres}
\subsection*{Ejercicio 2 \hfill 3,5 puntos} El 35\% de los habitantes de una ciudad lee el diario A, el 18\% el diario B, y el 8\% ambos diarios.

\begin{enumerate}[label=(\alph*)]

\item (1,75 puntos)  \textquestiondown{}Qué porcentaje de habitantes de esta ciudad no lee ninguno de los diarios?



\item (1,75 puntos)  Si se elige al azar un habitante de entre los que no leen el diario B, \textquestiondown{}cuál es la probabilidad de que lea el diario A?



\end{enumerate}

\clearpage
\section*{Bloque opcional (continuación)}
\subsection*{Ejercicio 3 \hfill 3,5 puntos} De los 800 alumnos matriculados en una asignatura, 320 son hombres y 480 mujeres. El 65\% de los hombres y el 75\% de las mujeres aprueban.

\begin{enumerate}[label=(\alph*)]

\item (2 puntos)  \textquestiondown{}Cuál es la probabilidad de que apruebe la asignatura?



\item (1,5 puntos)  Sabiendo que ha aprobado, \textquestiondown{}cuál es la probabilidad de que sea una mujer?



\end{enumerate}


\subsection*{Ejercicio 4 \hfill 3,5 puntos}
Una empresa dispone de dos secaderos, A y B. En el secadero B se curan el 40\% de los jamones. El 30\% de los jamones del secadero A y el 90\% de los del B son de tipo Reserva.

\begin{enumerate}[label=(\alph*)]

\item (2 puntos)  El jamón \textbf{no} es de Reserva.



\item (1,5 puntos)  Si el jamón es de Reserva, \textquestiondown{}cuál es la probabilidad de que proceda del secadero A?



\end{enumerate}


\end{document}
